O presente trabalho implementou os dois fluxos de processamento propostos sob as mesmas variáveis de ambiente, 
estrutura de rede e hardware, 
de tal forma que se distribuem da seguinte forma:

\begin{table}[htbp]
    \centering
    \caption{Variáveis de ambiente e parâmetros de configuração do sistema.}
    \label{tab:variaveis_ambiente}
    \begin{tabularx}{\textwidth}{lX}
        \toprule
        \textbf{Variável de Ambiente / Parâmetro} & \textbf{Configuração / Descrição} \\
        \midrule
        \textbf{Restrição Regulatória} & Ato nº 14.448 da ANATEL (Item 10.3, padrão AU915) \\
        \addlinespace
        \textbf{SF, DR, CR} & SF8, Data Rate DR6 (BW 500 kHz), Coding Rate 4/5 \\
        \addlinespace
        \textbf{Canal Fixo} & 915,9 MHz (Canal 64 do plano AU915, sem salto de frequência) \\
        \addlinespace
        \textbf{Potência de Transmissão} & +20 dBm (100 mW, fixa no módulo) \\
        \addlinespace
        \textbf{Tamanho de Imagens} & 160 $\times$ 120 pixels em tons de cinza ($\approx$ 4 KB comprimida em JPEG qualidade 90) \\
        \addlinespace
        \textbf{Versão do Python} & Python 3.10 (ambiente de treino) e Python 3.13 (ambiente de execução de scripts de testes) \\
        \addlinespace
        \textbf{Versão do TensorFlow} & TensorFlow 2.16.1 \\
        \addlinespace
        \textbf{Métodos de Ativação dos Nós} & \gls{abp} com chaves de sessão pré-configuradas \\
        \addlinespace
        \textbf{Portas de Aplicação (FPort 1 e 2)} & FPort 1 (blocos de imagem de até 175 bytes) e FPort 2 (leitura OCR de 9 bytes) \\
        \bottomrule
    \end{tabularx}
\end{table}

\paragraph{Variáveis utilizadas nos modelos de OCR implementados}
Durante o treinamento dos modelos de \gls{ocr} foram estabelecidas 
variáveis determinantes para a equiparação 
da comparação entre os modelos quantizados em INT8 e FP32.

\begin{table}[htbp]
    \centering
    \caption{Variáveis e hiperparâmetros definidos para o treinamento e quantização (FP32 e INT8)}
    \label{tab:hiperparametros}
    \begin{tabularx}{\textwidth}{l l X}
        \toprule
        \textbf{Variável / Parâmetro} & \textbf{Valor / Configuração} & \textbf{Função no Treinamento e Quantização} \\
        \midrule
        \texttt{args.epocas} & \texttt{60} (padrão) & Número máximo de épocas(sujeito ao \textit{Early Stopping}). \\
        \addlinespace
        \texttt{IMG\_SIZE} & \texttt{(28, 28)} & Dimensão espacial de entrada das imagens de dígitos redimensionadas. \\
        \addlinespace
        \texttt{SEED} & \texttt{42} & Semente de aleatoriedade global para garantir a reprodutibilidade. \\
        \addlinespace
        \texttt{BATCH\_SIZE} & \texttt{128} & Tamanho do lote (\textit{batch}) utilizado durante o treino e validação. \\
        \addlinespace
        \texttt{class\_weight} & Dicionário (máx. \texttt{3.0}) & Pesos calculados por classe para balancear a desproporção do dígito \texttt{'0'}. \\
        \addlinespace
        \texttt{augmentacao} & \texttt{Sequential} (5 camadas) & Aplica rotação, translação, contraste, ruído e \textit{downsample/upsample} simulando a câmera. \\
        \addlinespace
        \texttt{amostras\_repr} & Lote do dataset de treino & Amostras representativas usadas para calibrar a quantização INT8 no \texttt{TFLiteConverter}. \\
        \addlinespace
        \bottomrule
    \end{tabularx}
\end{table}

